\subsection{Strategi Pengujian}
\label{subsec:strategi-pengujian}

Strategi pengujian disusun berdasarkan kebutuhan sistem pada Tabel \ref{tab:kebutuhan-fungsional-tambahan-decmed} dan \ref{tab:kebutuhan-nonfungsional-tambahan-decmed}. Pengujian dibagi menjadi pengujian kebutuhan fungsional dan pengujian kebutuhan nonfungsional.

Pengujian kebutuhan fungsional diberi ID \texttt{PS-TA} dan dilakukan melalui skenario manual untuk memeriksa perilaku sistem dari sudut pandang aktor layanan kesehatan. Pengujian kebutuhan nonfungsional diberi ID \texttt{PU-TA} dan dilakukan melalui pengujian unit atau integrasi terhadap komponen token, segmentasi data, validasi akses, dan \textit{Server} PRE. Pengujian otomatis dilakukan dengan \texttt{cargo test} pada komponen yang relevan.

Kriteria keberhasilan pengujian ditentukan berdasarkan kesesuaian hasil skenario manual terhadap kebutuhan fungsional dan kelulusan pengujian otomatis terhadap kebutuhan nonfungsional.
